\documentclass[letterpaper,10pt,journal,twocolumn]{ieeeconf}

\usepackage{amsmath,amssymb,amsfonts,bm}
\usepackage{array}
\usepackage{booktabs}
\usepackage{cite}
\usepackage{url}
\usepackage{kotex}

\hyphenation{Kal-man Baye-sian state-space half-car quarter-car}

% 이 문서에서 실제로 사용하는 확률분포 및 연산자만 정의한다.
\newcommand{\Normal}{\mathcal{N}}
\newcommand{\Data}{\mathcal{D}}
\newcommand{\GP}{\mathcal{GP}}
\newcommand{\Bernoulli}{\operatorname{Bernoulli}}
\newcommand{\HalfNormal}{\operatorname{HalfNormal}}
\newcommand{\LogNormal}{\operatorname{LogNormal}}
\newcommand{\sigmoid}{\operatorname{sigmoid}}
\newcommand{\sat}{\operatorname{sat}}
\newcommand{\vecop}{\operatorname{vec}}
\newcommand{\iid}{\overset{\mathrm{iid}}{\sim}}

\allowdisplaybreaks[1]
\setlength{\emergencystretch}{1em}

\title{KF 기반 센서 재구성과 Bayesian Feature Modeling 설계}
\author{Jaewoong Han}

\begin{document}
\maketitle

\section{모델 체계와 기호}

이 문서에서 비교할 네 가지 KF 계열 모델은 다음과 같이 정리한다.

\begin{table*}[t]
\caption{비교할 네 가지 KF 계열 모델}
\label{tab:model-overview}
\centering
\scriptsize
\setlength{\tabcolsep}{3.5pt}
\renewcommand{\arraystretch}{1.18}
\begin{tabular}{c
                >{\raggedright\arraybackslash}p{0.19\textwidth}
                >{\raggedright\arraybackslash}p{0.25\textwidth}
                >{\raggedright\arraybackslash}p{0.20\textwidth}
                >{\raggedright\arraybackslash}p{0.23\textwidth}}
\toprule
번호 & 차량 모델 & Road/bump 처리 & 추정기 & 연구상 역할 \\
\midrule
1 & 축약 bounce--pitch body model
  & 명시적 road 추정 없음
  & Linear KF
  & 현재 코드에 가장 가까운 baseline \\
2 & 2-DOF quarter-car
  & Road height/rate를 상태에 증강
  & Augmented linear KF
  & 가장 단순한 road-estimation baseline \\
3 & 4-DOF linear half-car
  & 단일 spatial road profile과 front--rear delay
  & Augmented 또는 joint input--state KF + RTS smoother
  & 선형 주 후보 \\
4 & \texttt{vehicle\_model.py} nonlinear half-car
  & Bump profile/parameter를 latent state로 추정
  & EKF + smoother
  & 비선형 확장 후보 \\
\bottomrule
\end{tabular}
\end{table*}

기호는 다음처럼 구분한다.

\begin{itemize}

\item
  \(x_k\): 시각 \(k\)의 차량 동역학 상태
\item
  \(c_k\): 측정 가능한 제어·외생 입력, 예: wheel torque
\item
  \(r_k\): front/rear road height 또는 bump input
\item
  \(o_k\): IMU, wheel speed 등 실제 관측 센서
\item
  \(p\): 질량, 관성, 강성, 감쇠 등 차량 물리 파라미터
\item
  \(\tau_{ui}\): 사용자 \(u\), 구간 \(i\)에서 downstream hierarchical Bayesian model에 들어가는 feature
\item
  \(y_{ui}\): 기존 hierarchical Bayesian model의 outcome
\end{itemize}

중요하게도 \textbf{KF가 추정하는 state}, \textbf{road라는 unknown input}, \textbf{offline에서 맞추는 차량 파라미터}는 서로 다른 대상이다.

\section{큰 틀에서 KF의 기본 수식 설계}

\subsection{연속시간 차량 모델}

선형 차량 모델의 가장 일반적인 형태는 다음과 같다.

\[
\dot x(t)
=
A_c(p)x(t)
+B_c(p)c(t)
+E_c(p)r(t)
+w(t)
\]

\[
o(t)
=
C_c(p)x(t)
+D_c(p)c(t)
+J_c(p)r(t)
+v(t)
\]

여기서

\[
w(t)\sim\Normal(0,Q_c),
\qquad
v(t)\sim\Normal(0,R_c)
\]

이다. \(w\)는 process/model error, \(v\)는 sensor measurement error이다.

샘플링 간격 \(\Delta t\)로 이산화하면

\[
x_{k+1}
=
A_dx_k+B_dc_k+E_dr_k+w_k
\]

\[
o_k
=
C_dx_k+D_dc_k+J_dr_k+v_k
\]

이며, zero-order hold를 가정할 때

\[
\begin{aligned}
A_d &= e^{A_c\Delta t},\\
B_d &= \int_0^{\Delta t}e^{A_c\tau}B_c\,d\tau,\\
E_d &= \int_0^{\Delta t}e^{A_c\tau}E_c\,d\tau
\end{aligned}
\]

로 구성한다. 이는 Kalman의 선형 state-space filtering 틀에 해당한다 \cite{kalman1960}.

\subsection{\(A,B\)와 차량 파라미터의 offline 추정}

현재처럼 KF를 실행하기 전에 dynamics parameter를 calibration data로 맞출 수 있다.

\[
\hat p
=
\arg\min_p
\left[
\sum_{k=1}^{T}
\left\|
o_k-\hat o_k(p)
\right\|_{W}^{2}
+\lambda\mathcal R(p)
\right]
\]

그 후

\[
A_d^{*}=A_d(\hat p),
\qquad
B_d^{*}=B_d(\hat p)
\]

를 KF에서 사용한다.

가능하면 \(A,B\)의 원소를 완전히 자유롭게 맞추기보다

\[
\begin{aligned}
p=\{&
m_s,I_y,m_{uf},m_{ur},\\
&k_{sf},k_{sr},c_{sf},c_{sr},\\
&k_{tf},k_{tr},\ldots\}
\end{aligned}
\]

와 같은 물리 파라미터나, 안정성이 보장되는 modal parameter를 맞춘 뒤 \(A,B\)를 계산하는 편이 좋다. 그래야 음의 질량·감쇠나 불안정 pole 같은 비물리적 해를 방지할 수 있다.

\subsection{Kalman filter prediction--correction}

Road가 known input이면 \(r_k\)를 그대로 사용하고, road를 추정하지 않는 모델이면 해당 항을 제거한다.

\subsubsection{Prediction}

\[
\hat x_{k|k-1}
=
A_d\hat x_{k-1|k-1}
+B_dc_{k-1}
+E_d r_{k-1}
\]

\[
P_{k|k-1}
=
A_dP_{k-1|k-1}A_d^\top+Q
\]

\subsubsection{Innovation과 gain}

\[
e_k
=
o_k-
\left(
C_d\hat x_{k|k-1}
+D_dc_k
+J_dr_k
\right)
\]

\[
S_k
=
C_dP_{k|k-1}C_d^\top+R
\]

\[
K_k
=
P_{k|k-1}C_d^\top S_k^{-1}
\]

\subsubsection{Correction}

\[
\hat x_{k|k}
=
\hat x_{k|k-1}+K_ke_k
\]

수치적으로 안정적인 Joseph form은

\[
\begin{aligned}
P_{k|k}
={}&(I-K_kC_d)P_{k|k-1}(I-K_kC_d)^\top\\
&+K_kRK_k^\top
\end{aligned}
\]

이다.

\subsection{재구성 신호와 feature}

KF posterior로부터 측정 불가능한 센서 신호를 계산한다.

\[
\hat s_k
=
H_\tau\hat x_{k|k}
\]

예를 들어

\[
\hat s_k
=
\begin{bmatrix}
\widehat{\dot z_s}_k\\
\widehat{\dot\theta}_k
\end{bmatrix}
\]

가 bounce rate와 pitch rate가 된다. 한 window의 downstream feature는

\[
\hat\tau_{ui}^{KF}
=
\Phi
\left(
\hat s_{1:T},o_{1:T}
\right)
\]

로 계산한다. \(\Phi\)는 RMS, peak, standard deviation, derivative feature 등을 만드는 기존 feature extractor이다.

\subsection{Bump 통과가 끝난 뒤의 smoothing}

온라인 제어가 아니라 한 구간 전체가 수집된 뒤 feature를 복원하는 목적이라면 filtering만 쓰기보다 Rauch--Tung--Striebel smoother를 추가하는 것이 적절하다.

\[
G_k
=
P_{k|k}A_d^\top
\left(P_{k+1|k}\right)^{-1}
\]

\[
\hat x_{k|T}
=
\hat x_{k|k}
+G_k
\left(
\hat x_{k+1|T}-\hat x_{k+1|k}
\right)
\]

\[
\begin{aligned}
P_{k|T}
={}&P_{k|k}\\
&+G_k
\left(P_{k+1|T}-P_{k+1|k}\right)
G_k^\top
\end{aligned}
\]

Xue et al.도 half-car road reconstruction에서 augmented KF와 RTS smoothing을 결합한다 \cite{xue2020}.

\section{기존 Bayesian model에 추가하는 세 가지 설계}

세 방법은 같은 residual을 서로 다르게 표현한 것이 아니다.

\begin{itemize}

\item
  Feature residual: 최종 \(\tau\)의 복원오차를 보정
\item
  Dynamics residual: 매 time step의 상태전이 오차를 보정
\item
  Road bump estimation: vehicle error가 아니라 물리적인 미지 외란 \(r_k\)를 복원
\end{itemize}

초기 모델에서는 세 항을 모두 자유롭게 동시에 추정하면 안 된다. 특히 dynamics residual과 road input은 동일한 진동을 서로 대신 설명할 수 있어 식별성이 크게 떨어진다.

\subsection{Feature residual을 Bayesian하게 추정}

\subsubsection{Calibration data}

일부 calibration 구간에서 실제 target sensor를 측정할 수 있다고 가정한다.

\[
\hat\tau_{ui,j}^{KF}
=
\Phi_j
\left(
\hat x_{ui,1:T}
\right)
\]

\[
r_{ui,j}^{\tau}
=
\tau_{ui,j}^{true}
-\hat\tau_{ui,j}^{KF}
\]

\subsubsection{Bayesian residual regression}

\[
r_{ui,j}^{\tau}
=
q_{ui}^{\top}\beta_j
+b_{u,j}
+\epsilon_{ui,j}
\]

\[
\epsilon_{ui,j}
\iid
\Normal(0,\sigma_{\tau,j}^{2})
\]

\[
\begin{aligned}
\beta_j &\sim\Normal(0,\lambda_{\beta,j}^{2}I),\\
b_{u,j} &\sim\Normal(0,\omega_j^2),\\
\sigma_{\tau,j} &\sim\HalfNormal(s_j)
\end{aligned}
\]

여기서 \(\beta_j\)는 Beta 분포가 아니라, speed·torque·bump 조건 등이 residual 평균에 미치는 영향을 나타내는 \textbf{회귀계수 벡터}이다. 예를 들면

\[
q_{ui}
=
\begin{bmatrix}
1&
\bar v_{ui}&
T_{ui}^{peak}&
a_{x,ui}^{RMS}
\end{bmatrix}^{\top}
\]

로 둘 수 있다.

따라서 실측할 수 없는 feature의 조건부 평균을

\[
\mu_{ui,j}^{\tau}
=
\hat\tau_{ui,j}^{KF}
+q_{ui}^{\top}\beta_j
+b_{u,j}
\]

로 두면,

\[
\tau_{ui,j}^{mis}\mid\mu_{ui,j}^{\tau},\sigma_{\tau,j}
\sim
\Normal\!\left(\mu_{ui,j}^{\tau},\sigma_{\tau,j}^{2}\right)
\]

가 된다.

\(q_{ui}\)가 필요하지 않은 가장 단순한 시작 모델은

\[
\tau_{ui,j}^{mis}
\sim
\Normal
\left(
\hat\tau_{ui,j}^{KF}+b_j,
\sigma_{\tau,j}^{2}
\right)
\]

이다.

위 식은 KF posterior mean을 offset으로 쓰는 가장 단순한 형태이다. KF/EKF가 covariance까지 출력하도록 구현하면 더 완전하게

\[
x_{ui,1:T}^{(m)}
\sim
p_{KF}
\left(
x_{ui,1:T}\mid o_{ui,1:T}
\right)
\]

\[
\begin{aligned}
\tau_{ui,j}^{(m)}
={}&\Phi_j\!\left(x_{ui,1:T}^{(m)}\right)
+q_{ui}^{\top}\beta_j\\
&+b_{u,j}
+\epsilon_{ui,j}^{(m)}
\end{aligned}
\]

로 state-reconstruction uncertainty와 feature-residual uncertainty를 함께 전달할 수 있다. 따라서 현재 코드가 state mean만 반환한다면 \(P_{k|k}\) 또는 smoothed trajectory draw도 반환하도록 수정해야 한다.

\subsubsection{기존 hierarchical Bayesian model과의 결합}

기존 사용자별 coefficient 구조를 유지한다.

\[
\gamma_u
\sim
\Normal(\mu_\gamma,\Sigma_\gamma)
\]

관측·복원 feature를

\[
\widetilde\tau_{ui}
=
\begin{bmatrix}
\tau_{ui}^{obs}\\
\tau_{ui}^{mis}
\end{bmatrix}
\]

로 묶는다. 이진 outcome을 예로 들면

\[
y_{ui}\mid\widetilde\tau_{ui},\gamma_u
\sim
\Bernoulli\!\left[
\sigmoid\!\left(\gamma_u^\top\widetilde\tau_{ui}\right)
\right]
\]

최종 예측은 feature posterior를 적분한다.

\[
\begin{aligned}
p(y_{ui}\mid o_{ui},\Data)
={}&
\int p(y_{ui}\mid\tau_{ui},\gamma_u)\\
&\times p(\tau_{ui}^{mis}\mid o_{ui},\Data)
p(\gamma_u\mid\Data)
\,d\tau\,d\gamma
\end{aligned}
\]

즉 Normal 객체를 단순히 기존 \texttt{fit()}에 넣는 것이 아니라, posterior draw를 사용하거나 joint model 안에서 latent feature를 직접 샘플링한다.

이 설계와 통계적으로 가장 가까운 문헌은 Bayesian covariate-measurement-error와 latent-predictor model이다. Bartlett and Keogh는 true covariate를 latent하게 두고 outcome model까지 uncertainty를 전달하며, Schofield는 추정된 latent variable을 predictor로 쓸 때 measurement error를 함께 반영한다 \cite{bartlett2018,schofield2015}. 물리모델 출력 뒤의 discrepancy라는 관점에서는 Kennedy--O'Hagan calibration과 대응한다 \cite{kennedy2001}.

\subsection{Dynamics residual을 Bayesian하게 추정}

\(A^{*},B^{*}\) 또는 nonlinear dynamics parameter \(\hat p\)를 먼저 맞춘 뒤 다음과 같이 둔다.

\subsubsection{선형식}

\[
x_{k+1}
=
A^{*}x_k+B^{*}c_k+E^{*}r_k+d_k+w_k
\]

가장 단순한 systematic bias는

\[
d_k=b_x,
\qquad
b_x\sim\Normal(0,\Sigma_b)
\]

이다. 상태·입력 조건에 따라 error가 달라지면

\[
d_k
=
W\phi(x_k,c_k)
\]

\[
\vecop(W)
\sim
\Normal(0,\Lambda_W)
\]

로 둔다.

\subsubsection{Nonlinear식}

\[
x_{k+1}
=
F_d(x_k,c_k,r_k;\hat p)
+d(x_k,c_k)
+w_k
\]

\[
d(\cdot)\sim\GP(0,k_d)
\]

처럼 residual function에 Gaussian-process prior를 둘 수도 있다. BayesRace는 nominal vehicle model과 실제 transition 간 residual을 GP로 보정하며 \cite{jain2021}, GP state-space model은 latent state와 transition dynamics를 fully Bayesian하게 함께 추정한다 \cite{frigola2013}.

\subsubsection{무엇이 Bayesian 확장인가}

\[
w_k\sim\Normal(0,Q)
\]

를 고정된 \(Q\)와 함께 추가하는 것만으로는 이미 일반 KF와 같다. Bayesian dynamics residual이라고 부르려면 최소한

\[
p(b_x,W,Q\mid o_{1:T},c_{1:T})
\]

또는

\[
p(d(\cdot),Q\mid o_{1:T},c_{1:T})
\]

를 추정해야 한다.

100 Hz raw transition residual은 시간상관이 남을 가능성이 높으므로, IID가 맞지 않으면

\[
d_{k+1}
=
\rho_d d_k+\eta_k,
\qquad
\eta_k\sim\Normal(0,Q_d)
\]

와 같은 colored residual이 더 적절하다.

State posterior sample마다 feature를 계산하면

\[
x_{1:T}^{(m)}
\sim
p(x_{1:T}\mid o_{1:T})
\]

\[
\tau_{ui}^{(m)}
=
\Phi
\left(
x_{1:T}^{(m)}
\right)
\]

이므로 dynamics uncertainty가 downstream feature까지 전달된다.

\subsection{Road bump를 Bayesian하게 추정}

Road bump는 vehicle-model residual이 아니라 별도의 물리적 unknown input이다.

\[
x_{k+1}
=
F_d(x_k,c_k,r_{f,k},r_{r,k};p)
+w_k
\]

\[
o_k
=
h(x_k,c_k,r_{f,k},r_{r,k};p)
+v_k
\]

초기 bump 모델에서는 dynamics discrepancy \(d_k\)를 빼고, 고정 또는 사전 calibration된 vehicle model을 사용한다.

\subsubsection{하나의 공간 profile}

Front axle의 진행 위치를 \(s_k\)라고 하면

\[
s_{k+1}
=
s_k+v_{x,k}\Delta t
\]

\[
r_{f,k}
=
\rho(s_k)
\]

\[
r_{r,k}
=
\rho(s_k-L)
\]

이다. \(L=a+b\)는 wheelbase이다. 일정 속도일 때 이는

\[
r_r(t)
=
r_f
\left(
t-\frac{L}{v}
\right)
\]

와 같다. Zhu et al.도 front/rear road excitation을 wheelbase와 speed에 따른 delay로 연결한다 \cite{zhu2022}.

\subsubsection{형태가 알려진 speed bump: parametric Bayesian inverse problem}

Bump가 하나의 매끄러운 돌출부라고 가정하면

\[
\begin{aligned}
\rho(s;\alpha)
={}&
\frac{h_b}{2}
\left[
1-\cos\!\left(
\frac{2\pi(s-s_0)}{L_b}
\right)
\right]\\
&\times
\mathbf 1\!\left(0\le s-s_0\le L_b\right)
\end{aligned}
\]

\[
\alpha=
\begin{bmatrix}
h_b&L_b&s_0
\end{bmatrix}^{\top}
\]

로 둘 수 있다.

\[
\begin{aligned}
h_b &\sim\HalfNormal(s_h),\\
L_b &\sim\LogNormal(\mu_L,\sigma_L^2),\\
s_0 &\sim\Normal(\mu_s,\sigma_s^2)
\end{aligned}
\]

처럼 prior를 두고

\[
p(x_{0:T},\alpha\mid o_{1:T},c_{1:T})
\]

를 추정한다. 높이·폭·위치라는 소수의 파라미터만 추정하므로 현재처럼 vertical suspension sensor가 제한된 상황에서 가장 식별하기 쉽다.

\subsubsection{형태를 모르는 road: nonparametric spatial state}

공간 grid \(j\)에서

\[
\begin{bmatrix}
\rho_{j+1}\\
\rho'_{j+1}
\end{bmatrix}
=
\begin{bmatrix}
1&\Delta s\\
0&1
\end{bmatrix}
\begin{bmatrix}
\rho_j\\
\rho'_j
\end{bmatrix}
+\eta_j
\]

와 같은 smooth random-walk prior를 둘 수 있다. 이 경우 front와 rear가 같은 \(\rho\)를 서로 다른 위치에서 읽도록 buffer 또는 batch spatial smoother가 필요하다.

\subsubsection{대표 추정 방법}

이론적으로는 \textbf{joint state--unknown input estimation} 또는 \textbf{augmented-state nonlinear state-space estimation}이다.

\begin{itemize}

\item
  선형·온라인: augmented KF 또는 minimum-variance joint input--state estimator
\item
  비선형·온라인: augmented EKF/UKF
\item
  bump 통과 후 전체 복원: EKF/UKF + RTS-type smoother
\item
  posterior가 강하게 비Gaussian이면: particle smoother
\end{itemize}

Gillijns and De Moor는 state와 unknown input을 함께 minimum-variance unbiased하게 추정하는 재귀식을 제시한다 \cite{gillijns2007}. 차량 road-profile 사례로는 quarter-car augmented KF를 사용한 Doumiati et al.과, half-car에 AKF·regularization·RTS smoothing을 결합한 Xue et al., half-car joint input--state 추정을 사용한 Zhu et al.이 직접적인 근거이다 \cite{doumiati2011,xue2020,zhu2022}.

\subsection{세 Bayesian 설계의 최종 역할}

\begin{table*}[t]
\caption{세 Bayesian 설계의 역할}
\label{tab:bayesian-roles}
\centering
\scriptsize
\setlength{\tabcolsep}{4pt}
\renewcommand{\arraystretch}{1.18}
\begin{tabular}{
>{\raggedright\arraybackslash}p{0.18\textwidth}
>{\raggedright\arraybackslash}p{0.25\textwidth}
>{\raggedright\arraybackslash}p{0.29\textwidth}
>{\raggedright\arraybackslash}p{0.20\textwidth}}
\toprule
설계 & 보정·추정 대상 & 권장 위치 & 초기 연구에서의 역할 \\
\midrule
Feature residual
  & RMS·peak 등 최종 feature error
  & KF와 hierarchical outcome model 사이
  & 주 Bayesian feature model \\
Dynamics residual
  & 매 time-step vehicle transition error
  & KF/EKF process model 내부
  & 별도 ablation \\
Road bump
  & 물리적 외부 입력 \(\rho(s)\) 또는 bump parameter
  & KF/EKF state/input layer
  & Bump 실험에서 필수 \\
\bottomrule
\end{tabular}
\end{table*}

권장 조합은 다음과 같다.

\[
\boxed{
\begin{gathered}
\text{calibrated vehicle dynamics}
+\text{latent bump}\\
+\text{KF/EKF smoother}\\
+\text{Bayesian feature residual}\\
+\text{hierarchical outcome model}
\end{gathered}
}
\]

Dynamics residual과 bump를 동시에 자유롭게 추정하는 모델은 초기 주모델로 두지 않는다.

\section{KF 모델 4개 최종 내용}

\subsection{모델 1 --- 축약 bounce--pitch linear KF}

\subsubsection{근거 논문의 기본 차량 수식}

Half-car vertical dynamics는 일반적으로 다음 2차계 형태로 표현된다.

\[
M\ddot q+C\dot q+Kq
=
S_r r(t)+S_c c(t)
\]

\[
q=
\begin{bmatrix}
z_s&\theta&z_{uf}&z_{ur}
\end{bmatrix}^{\top}
\]

Zhu et al.도 차량의 mass, damping, stiffness matrix를 이용하여

\[
M\ddot Y+C\dot Y+KY=S_pF(t)
\]

로 쓰고, 이를

\[
\begin{aligned}
\dot X
={}&
\begin{bmatrix}
0&I\\
-M^{-1}K&-M^{-1}C
\end{bmatrix}X\\
&+
\begin{bmatrix}
0\\
M^{-1}S_p
\end{bmatrix}F
\end{aligned}
\]

로 변환한다 \cite[eqs.~(7)--(9)]{zhu2022}.

\begin{itemize}

\item
  차량 형태: front--rear를 포함하는 longitudinal half-car
\item
  원 논문의 online 추정 대상: vehicle response state와 unknown road force/input
\item
  원 논문의 고정 파라미터: \(M,C,K\)에 들어가는 질량·관성·강성·감쇠
\item
  현재 모델 1에서는 위 full half-car 중 low-frequency sprung-body bounce와 pitch mode만 남긴다.
\end{itemize}

\subsubsection{우리 수정 모델 수식}

현재 코드에 가까운 상태를

\[
\begin{aligned}
x_k
&=
\begin{bmatrix}
x_{b,k}^{\top}&d_k^{\top}
\end{bmatrix}^{\top},\\
x_{b,k}
&=
\begin{bmatrix}
z&\dot z&\theta&\dot\theta
\end{bmatrix}^{\top},\\
d_k
&=
\begin{bmatrix}
\eta_z&\eta_\theta&b_z&g_r
\end{bmatrix}^{\top}
\end{aligned}
\]

로 둔다.

연속시간 gray-box 식은 다음처럼 정리할 수 있다.

\[
\ddot z
=
-\omega_z^2z
-2\zeta_z\omega_z\dot z
+\kappa_a a_x
+\eta_z
\]

\[
\begin{aligned}
\ddot\theta
={}&
c_{z\theta}\dot z
-\omega_\theta^2\theta
-2\zeta_\theta\omega_\theta\dot\theta\\
&+g_a a_x
+g_TT_w
+g_{\Delta a}\Delta a_{fr}
+\eta_\theta
\end{aligned}
\]

Colored excitation과 bias는

\[
\eta_{z,k+1}
=
\rho_z\eta_{z,k}+\xi_{z,k}
\]

\[
\eta_{\theta,k+1}
=
\rho_\theta\eta_{\theta,k}+\xi_{\theta,k}
\]

\[
b_{z,k+1}=b_{z,k}+\xi_{b,k},
\qquad
g_{r,k+1}=g_{r,k}+\xi_{g,k}
\]

로 둔다.

관측식은

\[
o_{z,k}
=
\ddot z_k+b_{z,k}+\nu_{z,k}
\]

\[
o_{x,k}-\dot v_{x,k}
\simeq
g\theta_k+g\,g_{r,k}+b_x+\nu_{x,k}
\]

이다. IMU 좌표계·부호에 따라 \(g\theta\) 항의 부호는 실제 calibration으로 확정해야 한다.

\begin{itemize}

\item
  수정 이유:

  \begin{itemize}
  
  \item
    현재 센서에는 suspension deflection과 wheel vertical acceleration이 없으므로 full road-input model보다 관측 가능한 저차 body mode부터 안정적으로 복원한다.
  \item
    bump를 물리적인 road height로 분리하지 않고 \(\eta_z,\eta_\theta\)가 미모델링 excitation을 흡수한다.
  \item
    torque와 front--rear wheel-speed-derived acceleration difference를 pitch excitation으로 사용한다.
  \end{itemize}
\item
  사용하는 센서:

  \begin{itemize}
  
  \item
    vertical IMU acceleration: \(o_z\)
  \item
    longitudinal IMU acceleration: \(o_x\)
  \item
    네 바퀴 wheel speed: \(v_x\), \(\dot v_x\), \(\Delta a_{fr}\) 계산
  \item
    wheel/motor torque estimate: \(T_w\)
  \item
    roll rate: 이 2D 모델에서 추정하지 않고 측정값을 그대로 사용
  \end{itemize}
\item
  Online 추정 state:

  \begin{itemize}
  
  \item
    \(z,\dot z,\theta,\dot\theta\)
  \item
    colored excitation \(\eta_z,\eta_\theta\)
  \item
    vertical bias \(b_z\), grade-like state \(g_r\)
  \end{itemize}
\item
  Offline 최적화 대상:

  \begin{itemize}
  
  \item
    \(\omega_z,\zeta_z,\omega_\theta,\zeta_\theta\)
  \item
    \(c_{z\theta},\kappa_a,g_a,g_T,g_{\Delta a}\)
  \item
    \(Q,R\) 또는 그 diagonal scale
  \end{itemize}
\item
  최종 재구성:

  \begin{itemize}
  
  \item
    bounce rate: \(\widehat{\dot z}\)
  \item
    pitch rate: 부호정의에 따른 \(\pm\widehat{\dot\theta}\)
  \end{itemize}
\item
  한계:

  \begin{itemize}
  
  \item
    unsprung mass와 tire가 없으므로 bump height·width를 물리적으로 복원할 수 없다.
  \item
    이 모델에서 road는 상태가 아니라 process excitation이다.
  \end{itemize}
\end{itemize}

\subsection{모델 2 --- Quarter-car augmented KF}

\subsubsection{근거 논문 수식}

Doumiati et al.의 2-DOF quarter-car는 sprung mass \(m_s\)와 unsprung mass \(m_u\)로 구성된다.

\[
m_s\ddot z_s
=
-k_s(z_s-z_u)
-c_s(\dot z_s-\dot z_u)
\]

\[
m_u\ddot z_u
=
k_s(z_s-z_u)
+c_s(\dot z_s-\dot z_u)
-k_t(z_u-r)
-c_t(\dot z_u-\dot r)
\]

논문은 road를 unknown state에 포함하여

\[
x_k
=
\begin{bmatrix}
z_s&
\dot z_s&
z_u&
\dot z_u&
r&
\dot r
\end{bmatrix}^{\top}
\]

\[
x_{k+1}=A_kx_k+w_k
\]

로 둔다. Road prior는

\[
\ddot r=0
\]

인 constant-slope/random-walk 형태이다. 논문의 observation은

\[
o_k
=
\begin{bmatrix}
z_s-z_u&
z_s&
\ddot z_s
\end{bmatrix}^{\top}
\]

이다 \cite[eqs.~(1)--(4)]{doumiati2011}.

Chaabane et al.도 quarter-car에서 road input을 상태에 붙여

\[
x_{k+1}^{a}
=
\begin{bmatrix}
A_d&B_d\\
0&I
\end{bmatrix}
x_k^{a}
+
\begin{bmatrix}
w_k\\
\eta_k
\end{bmatrix}
\]

\[
o_k=[C\;D]x_k^a+v_k
\]

인 augmented KF를 사용한다 \cite[eq.~(3.5)]{chaabane2019}.

\begin{itemize}

\item
  차량 형태: 차체 1/4과 한 wheel/axle을 나타내는 2-DOF quarter-car
\item
  원 논문의 online 추정 대상:

  \begin{itemize}
  
  \item
    \(z_s,\dot z_s,z_u,\dot z_u\)
  \item
    road height \(r\)와 road rate \(\dot r\)
  \end{itemize}
\item
  원 논문의 고정 파라미터:

  \begin{itemize}
  
  \item
    \(m_s,m_u,k_s,c_s,k_t,c_t\)
  \end{itemize}
\item
  원 논문의 센서:

  \begin{itemize}
  
  \item
    vertical acceleration
  \item
    suspension deflection
  \item
    적분으로 얻은 body displacement
  \end{itemize}
\item
  pitch와 roll은 모델에 포함되지 않는다.
\end{itemize}

\subsubsection{우리 수정 모델 수식}

현재 센서만 사용할 때 vertical observation은

\[
o_k
=
a_{z,k}^{IMU}
\]

로 단순화한다. 여기서

\[
a_{z,k}^{IMU}
=
\ddot z_{s,k}+b_{z,k}+\nu_{z,k}
\]

이다. Wheel speed로 얻은 \(v_x\)는 quarter-car vertical KF의 correction에 직접 넣기보다

\[
s_{k+1}=s_k+\bar v_{w,k}\Delta t
\]

로 road를 시간축에서 공간축으로 바꾸는 보조 입력으로 사용한다.

Speed bump의 형태가 대략 알려져 있다면 자유로운 \((r,\dot r)\) random walk보다

\[
r_k=\rho(s_k;\alpha),
\qquad
\alpha=
\begin{bmatrix}
h_b&L_b&s_0
\end{bmatrix}^{\top}
\]

로 제한하는 것이 현재 센서 구성에서는 더 적절하다.

\begin{itemize}

\item
  수정 이유:

  \begin{itemize}
  
  \item
    논문이 사용하는 suspension deflection sensor가 현재 데이터에는 없으므로 동일한 observability를 주장할 수 없다.
  \item
    자유 road state는 body vertical acceleration 하나와 강하게 confound되므로, known bump shape 또는 informative prior가 필요하다.
  \item
    quarter-car는 pitch를 표현하지 못하므로 pitch-rate reconstruction의 최종 모델이 아니라 road-estimation ablation으로 사용한다.
  \end{itemize}
\item
  사용하는 센서:

  \begin{itemize}
  
  \item
    vertical IMU acceleration
  \item
    네 wheel speed 평균으로 얻은 \(v_x\)와 이동거리 \(s\)
  \item
    torque는 순수 quarter-car vertical 식에는 직접 사용하지 않음
  \end{itemize}
\item
  Online 추정 대상:

  \begin{itemize}
  
  \item
    \(z_s,\dot z_s,z_u,\dot z_u\)
  \item
    제한된 road state 또는 bump parameter
  \item
    필요하면 \(b_z\)
  \end{itemize}
\item
  Offline 고정·calibration 대상:

  \begin{itemize}
  
  \item
    \(m_s,m_u,k_s,c_s,k_t,c_t,Q,R\)
  \end{itemize}
\item
  한계:

  \begin{itemize}
  
  \item
    현재 센서만으로 arbitrary road profile amplitude를 안정적으로 식별하기 어렵다.
  \item
    bounce 검증용 baseline에는 유용하지만 pitch target에는 구조적으로 부족하다.
  \end{itemize}
\end{itemize}

\subsection{모델 3 --- Linear half-car joint state--road estimation}

\subsubsection{근거 논문 수식}

Zhu et al.은 half-car 계열 차량의 동역학을

\[
M\ddot Y+C\dot Y+KY=S_pF(t)
\]

로 두고

\[
X_{k+1}
=
A_dX_k+B_dF_k+w_k
\]

\[
Z_k
=
GX_k+JF_k+v_k
\]

로 이산화한다. Acceleration, velocity, displacement selection matrix를 \(S_a,S_v,S_d\)라 하면

\[
G
=
\begin{bmatrix}
S_d-S_aM^{-1}K&
S_v-S_aM^{-1}C
\end{bmatrix}
\]

\[
J=S_aM^{-1}S_p
\]

이다 \cite[eqs.~(7)--(13)]{zhu2022}.

해당 논문의 unknown-input update는

\[
\widetilde R_k
=
GP_{k|k-1}G^\top+R
\]

\[
M_k
=
\left(
J^\top\widetilde R_k^{-1}J
\right)^{-1}
J^\top\widetilde R_k^{-1}
\]

\[
\hat F_{k|k}
=
M_k
\left(
Z_k-G\hat X_{k|k-1}
\right)
\]

로 road input을 먼저 구한 뒤 state를 update한다.

Xue et al.은 bounce와 pitch를 포함하는 half-car를 사용하고, front와 rear에서 추정한 동일 road profile의 일치도를 목적함수로 삼아 vehicle parameter를 최적화한 뒤 AKF, regularization, RTS smoothing으로 road를 복원한다 \cite{xue2020}.

\begin{itemize}

\item
  차량 형태:

  \begin{itemize}
  
  \item
    Zhu: equipment를 포함한 상세 2D half-vehicle
  \item
    Xue: bounce--pitch와 front/rear unsprung motion을 포함한 half-car
  \end{itemize}
\item
  원 논문의 online 추정 대상:

  \begin{itemize}
  
  \item
    body·axle state
  \item
    front/rear road input
  \end{itemize}
\item
  Xue의 offline identification:

  \begin{itemize}
  
  \item
    front/rear road-profile consistency를 이용한 vehicle mechanical parameter 최적화
  \end{itemize}
\item
  원 논문의 핵심 제약:

  \begin{itemize}
  
  \item
    front와 rear wheel은 같은 road를 wheelbase/speed만큼 시간차를 두고 통과
  \end{itemize}
\end{itemize}

\subsubsection{우리 수정 차량 수식}

현재 목적에는 standard 4-DOF passenger half-car로 단순화한다.

\[
\delta_f=z_s+a\theta-z_{uf},
\qquad
\delta_r=z_s-b\theta-z_{ur}
\]

\[
F_{sf}
=
k_{sf}\delta_f
+c_{sf}\dot\delta_f
\]

\[
F_{sr}
=
k_{sr}\delta_r
+c_{sr}\dot\delta_r
\]

\[
F_{tf}
=
k_{tf}(z_{uf}-r_f),
\qquad
F_{tr}
=
k_{tr}(z_{ur}-r_r)
\]

정적 평형을 원점으로 두고 upward-positive sign convention을 쓰면

\[
m_s\ddot z_s
=
-F_{sf}-F_{sr}
\]

\[
I_y\ddot\theta
=
-aF_{sf}
+bF_{sr}
+M_T
\]

\[
m_{uf}\ddot z_{uf}
=
F_{sf}-F_{tf}
\]

\[
m_{ur}\ddot z_{ur}
=
F_{sr}-F_{tr}
\]

Torque-induced pitch moment는 가장 단순하게

\[
M_T=\lambda_TT_w
\]

로 두거나, longitudinal tire force를 계산할 수 있으면

\[
F_x=\frac{T_w}{R_{\mathrm{eff}}},
\qquad
M_T=-h_{cg}F_x
\]

로 물리 경로를 사용한다.

상태는

\[
\begin{aligned}
q&=
\begin{bmatrix}
z_s&\theta&z_{uf}&z_{ur}
\end{bmatrix}^{\top},\\
x&=
\begin{bmatrix}
q^\top&\dot q^\top
\end{bmatrix}^{\top}
\end{aligned}
\]

이다.

Road는 두 개의 독립 random walk가 아니라

\[
r_{f,k}=\rho(s_k),
\qquad
r_{r,k}=\rho(s_k-L)
\]

로 묶는다.

현재 센서에 맞춘 observation은

\[
a_{z,k}^{IMU}
=
\ddot z_{s,k}
+x_I\ddot\theta_k
+b_{z,k}
+\nu_{z,k}
\]

\[
a_{x,k}^{IMU}-\dot v_{x,k}
\simeq
g\theta_k+b_{x,k}+\nu_{x,k}
\]

\[
v_{w,k}^{f/r}
\simeq
v_{x,k}+\nu_{v,k}^{f/r}
\]

로 둔다.

\begin{itemize}

\item
  수정 이유:

  \begin{itemize}
  
  \item
    원 Zhu 모델의 equipment DOF는 일반 passenger vehicle 센서 재구성에 불필요하므로 제거한다.
  \item
    현재 센서에는 suspension displacement와 axle vertical acceleration이 없으므로, \(J\)가 두 road input에 대해 full column rank라고 보장할 수 없다.
  \item
    따라서 instant direct-input update보다 하나의 spatial road profile, front--rear delay, informative prior, full-window smoothing을 함께 사용한다.
  \item
    longitudinal torque가 pitch에 주는 영향을 명시한다.
  \end{itemize}
\item
  사용하는 센서:

  \begin{itemize}
  
  \item
    vertical IMU acceleration
  \item
    longitudinal IMU acceleration
  \item
    front/rear axle별 평균 wheel speed
  \item
    wheel/motor torque estimate
  \item
    roll rate는 2D half-car 밖이므로 측정값 pass-through
  \end{itemize}
\item
  Online 또는 smoothed 추정 대상:

  \begin{itemize}
  
  \item
    \(z_s,\theta,z_{uf},z_{ur}\)와 각 속도
  \item
    bounce rate \(\dot z_s\), pitch rate \(\dot\theta\)
  \item
    하나의 road profile \(\rho(s)\) 또는 bump parameter \(\alpha\)
  \item
    IMU bias \(b_z,b_x\)
  \end{itemize}
\item
  Offline 최적화 대상:

  \begin{itemize}
  
  \item
    우선 \(k_{sf},k_{sr},c_{sf},c_{sr},k_{tf},k_{tr},\lambda_T,Q,R\)
  \item
    질량·관성은 차량 제원값을 우선 고정
  \item
    모든 질량·강성·감쇠와 bump amplitude를 동시에 자유롭게 맞추지 않음
  \end{itemize}
\item
  장점:

  \begin{itemize}
  
  \item
    모델 2와 달리 bounce와 pitch를 동시에 표현한다.
  \item
    front--rear delay가 bump의 식별성을 높인다.
  \end{itemize}
\item
  한계:

  \begin{itemize}
  
  \item
    current sensor set에서는 road height와 suspension parameter가 여전히 confound될 수 있으므로 observability/Fisher-information 검사가 필요하다.
  \end{itemize}
\end{itemize}

\subsection{모델 4 --- \texttt{vehicle\_model.py} 기반 nonlinear EKF}

\subsubsection{EKF의 기본 원리}

Nonlinear state-space model은

\[
x_{k+1}
=
F_d(x_k,c_k)+w_k
\]

\[
o_k=h(x_k,c_k)+v_k
\]

로 쓴다. EKF는 현재 추정점에서 1차 Taylor linearization을 수행한다. 비선형 차량 상태추정에서는 UKF도 대표적인 대안이며 \cite{antonov2011}, nonlinear half-vehicle multibody model과 augmented/constrained EKF를 결합한 직접적인 사례도 있다 \cite{debarbouille2022}.

Jacobian은

\[
F_k
=
\left.
\frac{\partial F_d}{\partial x}
\right|_{\hat x_{k|k},c_k}
\]

\[
H_k
=
\left.
\frac{\partial h}{\partial x}
\right|_{\hat x_{k|k-1},c_k}
\]

이다.

Prediction은

\[
\hat x_{k+1|k}
=
F_d(\hat x_{k|k},c_k)
\]

\[
P_{k+1|k}
=
F_kP_{k|k}F_k^\top+Q
\]

이고, correction은

\[
e_{k+1}
=
o_{k+1}-h(\hat x_{k+1|k},c_{k+1})
\]

\[
S_{k+1}
=
H_{k+1}P_{k+1|k}H_{k+1}^\top+R
\]

\[
K_{k+1}
=
P_{k+1|k}H_{k+1}^\top S_{k+1}^{-1}
\]

\[
\hat x_{k+1|k+1}
=
\hat x_{k+1|k}+K_{k+1}e_{k+1}
\]

이다.

\subsubsection{\texttt{vehicle\_model.py}의 원래 형태}

앞서 확인한 코드의 상태는

\[
\begin{aligned}
v&=
\begin{bmatrix}
\dot z_s&\dot\theta&\dot z_{uf}&\dot z_{ur}&v_x
\end{bmatrix}^{\top},\\
q&=
\begin{bmatrix}
z_s&\theta&z_{uf}&z_{ur}&s
\end{bmatrix}^{\top},\\
X&=
\begin{bmatrix}
v^\top&q^\top
\end{bmatrix}^{\top}
\end{aligned}
\]

이며

\[
\dot X
=
f_{\mathrm{veh}}
\left(
X,T_w,r_f,r_r;p
\right)
\]

형태이다.

\begin{itemize}

\item
  차량 형태: bounce, pitch, front/rear unsprung motion, longitudinal speed를 결합한 nonlinear 4-DOF half-car
\item
  원 코드의 known input:

  \begin{itemize}
  
  \item
    wheel/motor torque
  \item
    front/rear road profile \(r_f,r_r\)
  \end{itemize}
\item
  원 코드의 dynamic state:

  \begin{itemize}
  
  \item
    sprung/unsprung vertical velocities and displacements
  \item
    pitch rate and pitch angle
  \item
    longitudinal speed and traveled position
  \end{itemize}
\item
  원 코드에서 road는 추정 대상이 아니라 주어진 입력이었다.
\end{itemize}

\subsubsection{우리 수정 nonlinear dynamics}

Suspension force를

\[
F_{sf}
=
f_{sf}
\left(
\delta_f,\dot\delta_f;p
\right)
\]

\[
F_{sr}
=
f_{sr}
\left(
\delta_r,\dot\delta_r;p
\right)
\]

로 두면 linear spring--damper, nonlinear stiffness, saturation을 모두 포함할 수 있다.

\[
F_{tf}
=
f_{tf}(z_{uf}-r_f),
\qquad
F_{tr}
=
f_{tr}(z_{ur}-r_r)
\]

Torque와 longitudinal force는

\[
F_x
=
\sat
\left(
\frac{T_w}{R_{\mathrm{eff}}}
\right)
\]

\[
\dot v_x
=
\frac{
F_x-F_{\mathrm{roll}}(v_x)-F_{\mathrm{aero}}(v_x)
}{m_{\mathrm{tot}}}
\]

로 둔다.

Vertical/pitch dynamics는

\[
\ddot z_s
=
\frac{-F_{sf}-F_{sr}}{m_s}
\]

\[
\ddot\theta
=
\frac{-aF_{sf}+bF_{sr}-h_{cg}F_x}{I_y}
\]

\[
\ddot z_{uf}
=
\frac{F_{sf}-F_{tf}}{m_{uf}}
\]

\[
\ddot z_{ur}
=
\frac{F_{sr}-F_{tr}}{m_{ur}}
\]

로 쓸 수 있다. 실제 구현에서는 \texttt{vehicle\_model.py}가 사용하는 좌표·pitch 부호와 위 식을 하나로 통일해야 한다.

\subsubsection{Known road를 latent bump로 변경}

원래 입력이던 \(r_f,r_r\)를

\[
r_{f,k}
=
\rho(s_k;\alpha)
\]

\[
r_{r,k}
=
\rho(s_k-L;\alpha)
\]

로 교체한다.

Parametric bump이면 augmented state를

\[
X_k^a
=
\begin{bmatrix}
X_k^\top&\alpha^\top&b_z&b_x
\end{bmatrix}^{\top}
\]

로 두고

\[
\alpha_{k+1}
=
\alpha_k+\eta_{\alpha,k}
\]

로 천천히 변하거나 거의 고정된 parameter처럼 추정한다. 실제로 하나의 bump만 통과한다면 \(\eta_{\alpha,k}\)를 매우 작게 두고, filter 후 batch smoother 또는 MAP refinement를 수행하는 편이 안정적이다.

Actuator lag가 중요하면 실제 wheel torque도 state로 둔다.

\[
T_{w,k+1}
=
a_TT_{w,k}
+b_TT_{\mathrm{cmd},k}
+\xi_{T,k}
\]

\[
T_{\mathrm{est},k}
=
T_{w,k}+\nu_{T,k}
\]

\subsubsection{우리 sensor observation model}

\[
a_{z,k}^{IMU}
=
\ddot z_{s,k}
+x_I\ddot\theta_k
+b_{z,k}
+\nu_{z,k}
\]

\[
a_{x,k}^{IMU}
=
\dot v_{x,k}
+g\sin\theta_k
+b_{x,k}
+\nu_{x,k}
\]

작은 pitch에서는 \(g\sin\theta\simeq g\theta\)이다.

\[
v_{w,k}^{f}
\simeq
v_{x,k}+\nu_{vf,k}
\]

\[
v_{w,k}^{r}
\simeq
v_{x,k}+\nu_{vr,k}
\]

Slip이 큰 구간에서는 wheel speed를 곧바로 \(v_x\)로 쓰지 말고 covariance를 키우거나 slip state를 추가해야 한다.

\begin{itemize}

\item
  수정 이유:

  \begin{itemize}
  
  \item
    원 코드의 known road input을 bump latent variable로 전환해야 한다.
  \item
    position \(s\)는 road가 known time series일 때는 제거 가능하지만, 공간 bump \(\rho(s)\)를 추정할 때는 반드시 유지하는 것이 유리하다.
  \item
    torque command와 실제 torque estimate의 차이가 크면 actuator state를 추가한다.
  \item
    accelerometer bias를 state로 추가하여 gravity projection과 low-frequency drift를 분리한다.
  \item
    \(h_b,L_b>0\) 제약은 \(\log h_b,\log L_b\)를 state로 두어 만족시키고, bump 시작·끝의 indicator는 smooth gate로 근사하거나 filter 후 batch MAP에서 다룬다.
  \item
    saturation의 hard corner는 Jacobian이 불연속이므로 smooth saturation을 쓰거나, 비선형성이 강하면 UKF로 교체한다.
  \end{itemize}
\item
  사용하는 센서:

  \begin{itemize}
  
  \item
    vertical IMU acceleration
  \item
    longitudinal IMU acceleration
  \item
    four-wheel speed; front/rear average와 slip consistency 계산
  \item
    motor/wheel torque command 및 estimate
  \item
    roll rate는 2D 모델에 포함하지 않고 measured channel을 그대로 사용
  \end{itemize}
\item
  EKF/smoother 추정 대상:

  \begin{itemize}
  
  \item
    \(\dot z_s,\dot\theta,\dot z_{uf},\dot z_{ur},v_x\)
  \item
    \(z_s,\theta,z_{uf},z_{ur},s\)
  \item
    bump parameter \((h_b,L_b,s_0)\) 또는 제한된 spatial profile
  \item
    \(b_z,b_x\)와 필요 시 actual torque \(T_w\)
  \end{itemize}
\item
  Offline 고정·calibration 대상:

  \begin{itemize}
  
  \item
    질량, 관성, wheelbase, CG height
  \item
    suspension/tire stiffness and damping
  \item
    rolling/aero resistance coefficient
  \item
    \(Q,R\)와 초기 covariance
  \end{itemize}
\item
  추정하지 말아야 할 초기 항목:

  \begin{itemize}
  
  \item
    자유로운 vehicle dynamics residual과 자유로운 road profile을 동시에 추정하지 않음
  \item
    모든 물리 파라미터와 bump height를 동시에 online state로 두지 않음
  \end{itemize}
\item
  검증:

  \begin{itemize}
  
  \item
    calibration 구간에서 실제 bounce/pitch-rate sensor는 EKF observation에 넣지 않고 reconstruction target/validation으로 남긴다.
  \item
    normalized innovation squared, state covariance consistency, bump parameter credible interval을 함께 평가한다.
  \end{itemize}
\end{itemize}

\section{최종 연구 구성}

\subsection{권장 비교 실험}

\begin{enumerate}
\def\labelenumi{\arabic{enumi}.}

\item
  \textbf{Model 1: reduced linear KF}

  \begin{itemize}
  
  \item
    현재 baseline
  \item
    point feature만 downstream model에 입력
  \end{itemize}
\item
  \textbf{Model 1 + Bayesian feature residual}

  \begin{itemize}
  
  \item
    dynamics는 그대로 유지
  \item
    교수님의 feature-level Bayesian modeling 효과만 분리해 검증
  \end{itemize}
\item
  \textbf{Model 3: linear half-car + latent bump + RTS smoother}

  \begin{itemize}
  
  \item
    bump를 물리적 unknown input으로 추정
  \item
    bounce/pitch feature를 posterior sample로 생성
  \end{itemize}
\item
  \textbf{Model 4: nonlinear \texttt{vehicle\_model.py} EKF + latent bump + smoother}

  \begin{itemize}
  
  \item
    nonlinear dynamics의 추가 이득 검증
  \item
    Model 3 대비 성능이 실제로 좋아지는지 비교
  \end{itemize}
\item
  \textbf{Dynamics residual model}

  \begin{itemize}
  
  \item
    별도 ablation
  \item
    bump model과 동시에 자유롭게 두지 않고, known/controlled-road calibration에서 먼저 검증
  \end{itemize}
\end{enumerate}

\subsection{최종 주모델}

현재 센서 제약과 bump 통과 상황을 함께 고려하면 가장 방어하기 쉬운 주모델은

\[
\boxed{
\begin{gathered}
\text{Model 3 또는 4}
+\text{latent spatial bump}\\
+\text{front--rear wheelbase constraint}
+\text{smoother}\\
+\text{Bayesian feature residual}
\end{gathered}
}
\]

이다.

Model 3은 선형성과 식별성 때문에 주 분석에 더 안전하고, Model 4는 nonlinear extension 또는 성능 비교모델로 두는 것이 합리적이다. \texttt{vehicle\_model.py}의 nonlinear EKF가 충분히 안정적이고 실제 hold-out data에서 Model 3보다 좋아질 때 Model 4를 최종 주모델로 승격한다.

\subsection{반드시 명시할 한계}

\begin{itemize}

\item
  현재 sensor set에는 suspension deflection과 wheel vertical acceleration이 없으므로, 기존 road-profile 논문과 동일한 관측가능성을 자동으로 얻지 못한다.
\item
  Wheel speed는 bump의 공간 위치와 front--rear delay를 알려주지만 vertical amplitude를 직접 측정하지 않는다.
\item
  Bump height, tire stiffness, suspension stiffness를 모두 자유롭게 추정하면 scale confounding이 생긴다.
\item
  2D half-car는 roll을 추정하지 못하므로 roll은 측정 channel을 유지하거나 별도의 full-car/roll model이 필요하다.
\item
  한 번도 target sensor를 측정하지 않은 경우 feature residual model의 bias와 variance를 데이터만으로 식별할 수 없다. 적어도 calibration 구간의 target sensor 또는 강한 물리 prior가 필요하다.
\end{itemize}

\begin{thebibliography}{99}

\bibitem{kalman1960}
R.~E. Kalman, ``A new approach to linear filtering and prediction
problems,'' \emph{Journal of Basic Engineering}, vol.~82, no.~1,
pp.~35--45, 1960, doi: 10.1115/1.3662552.

\bibitem{xue2020}
K. Xue, T. Nagayama, and B. Zhao, ``Road profile estimation and
half-car model identification through the automated processing of
smartphone data,'' \emph{Mechanical Systems and Signal Processing},
vol.~142, Art.~no.~106722, 2020,
doi: 10.1016/j.ymssp.2020.106722.

\bibitem{bartlett2018}
J.~W. Bartlett and R.~H. Keogh, ``Bayesian correction for covariate
measurement error: A frequentist evaluation and comparison with
regression calibration,'' \emph{Statistical Methods in Medical Research},
vol.~27, no.~6, pp.~1695--1708, 2018,
doi: 10.1177/0962280216667764.

\bibitem{schofield2015}
L.~S. Schofield, ``Correcting for measurement error in latent variables
used as predictors,'' \emph{The Annals of Applied Statistics}, vol.~9,
no.~4, pp.~2133--2152, 2015, doi: 10.1214/15-AOAS877.

\bibitem{kennedy2001}
M.~C. Kennedy and A. O'Hagan, ``Bayesian calibration of computer
models,'' \emph{Journal of the Royal Statistical Society: Series B},
vol.~63, no.~3, pp.~425--464, 2001,
doi: 10.1111/1467-9868.00294.

\bibitem{jain2021}
A. Jain, M. O'Kelly, P. Chaudhari, and M. Morari, ``BayesRace:
Learning to race autonomously using prior experience,'' in
\emph{Proc. Conf. Robot Learning}, PMLR, vol.~155,
pp.~1918--1929, 2021.

\bibitem{frigola2013}
R. Frigola, F. Lindsten, T.~B. Sch\"on, and C.~E. Rasmussen,
``Bayesian inference and learning in Gaussian process state-space
models with particle MCMC,'' in \emph{Advances in Neural Information
Processing Systems 26}, 2013.

\bibitem{zhu2022}
J. Zhu, X. Chang, X. Zhang, Y. Su, and X. Long, ``A novel method for
the reconstruction of road profiles from measured vehicle responses
based on the Kalman filter method,'' \emph{Computer Modeling in
Engineering \& Sciences}, vol.~130, no.~3, 2022,
doi: 10.32604/cmes.2022.019140.

\bibitem{gillijns2007}
S. Gillijns and B. De Moor, ``Unbiased minimum-variance input and
state estimation for linear discrete-time systems,'' \emph{Automatica},
vol.~43, no.~1, pp.~111--116, 2007,
doi: 10.1016/j.automatica.2006.08.002.

\bibitem{doumiati2011}
M. Doumiati, A. Victorino, A. Charara, and D. Lechner,
``Estimation of road profile for vehicle dynamics motion:
Experimental validation,'' in \emph{Proc. American Control Conference},
pp.~5237--5242, 2011.

\bibitem{chaabane2019}
M.~M. Chaabane, D.~B. Hassen, M.~S. Abbes, S.~C. Baslamisli,
F. Chaari, and M. Haddar, ``Road profile identification using
estimation techniques: Comparison between independent component
analysis and Kalman filter,'' \emph{Journal of Theoretical and Applied
Mechanics}, vol.~57, no.~2, pp.~397--409, 2019,
doi: 10.15632/jtam-pl/104592.

\bibitem{antonov2011}
S. Antonov, A. Fehn, and A. Kugi, ``Unscented Kalman filter for
vehicle state estimation,'' \emph{Vehicle System Dynamics}, vol.~49,
no.~9, pp.~1497--1520, 2011,
doi: 10.1080/00423114.2010.527994.

\bibitem{debarbouille2022}
A. D\'ebarbouill\'e, F. Renaud, Z. Dimitrijevic, D. Chojnacki,
L. Rota, and J.-L. Dion, ``Wheel forces estimation with an augmented
and constrained extended Kalman filter applied on a nonlinear
multi-body model of a half vehicle,'' \emph{Procedia Structural
Integrity}, vol.~38, pp.~342--351, 2022,
doi: 10.1016/j.prostr.2022.03.035.

\end{thebibliography}

\end{document}
